\section{Shortening families in a fixed degree}\label{sec:shortening}

We first prove a filling estimate whose constant may depend on the homological
degree.  The geometric argument works on every compact simply connected
Riemannian manifold.  Its two ingredients are continuous arclength
parametrization and a shortening procedure with only finitely many simultaneous
changes.  Throughout this section, chains have coefficients in $\kk$;
for the manifold $M$ under consideration, write
$\Lambda=H^1(S^1,M)$ and $\mathcal A_a=\{\gamma\in\Lambda:\ell(\gamma)<a\}$.
Block, Manin and Weinberger established fixed-degree additive filling
bounds for based and free Lipschitz loop spaces
\cite[Theorem~5.1]{BlockManinWeinberger2025}. We give a direct proof in
the strong $H^1$ topology for the length filtration used here.

\subsection{Arclength parametrization in \texorpdfstring{$H^1$}{H1}}

\begin{lemma}[Continuous arclength homotopy]\label{lem:arclength}
Let $M$ be a compact Riemannian manifold.  There is a jointly continuous map
\[
 \Psi:[0,1]\times H^1([0,1],M)\longrightarrow H^1([0,1],M)
\]
in the strong $H^1$ topology such that $\Psi_0$ is the identity,
$\Psi_1\gamma$ has constant speed, and every $\Psi_u$ preserves length and
endpoints.  Constant paths and constant-speed paths are fixed.  Moreover,
\begin{equation}\label{eq:shortening-energy}
 E(\Psi_u\gamma)
 \le (1-u)E(\gamma)+u\ell(\gamma)^2
 \le E(\gamma).
\end{equation}
The homotopy commutes with reversal.  In particular, on loops it preserves
evaluation at zero and the constant-loop subspace.
\end{lemma}

\begin{proof}
Write $v(t)=|\dot\gamma(t)|$ and $L=\int_0^1v(t)\,\dd t$.
When $L>0$, put
\[
 \sigma(t)=\frac1L\int_0^t v(s)\,\dd s,
 \qquad \phi_u(t)=(1-u)t+u\sigma(t).
\]
For $u<1$, the absolutely continuous map $\phi_u$ is an increasing
homeomorphism with Lipschitz inverse.  Set
$\Psi_u\gamma=\gamma\circ\phi_u^{-1}$.
At $u=1$ the inverse need not exist.  Instead, define
$\Psi_1\gamma(s)$ to be the value of $\gamma$ on the fiber
$\sigma^{-1}(s)$.  Each fiber is an interval on which $v=0$ almost
everywhere, so absolute continuity makes this definition unambiguous.
For $L=0$ leave the constant path fixed for every $u$.

For $s_1<s_2$, choose ordered preimages $t_1<t_2$ under $\sigma$.  Then
\[
 d\bigl(\Psi_1\gamma(s_1),\Psi_1\gamma(s_2)\bigr)
 \le\int_{t_1}^{t_2}v(t)\,\dd t=L(s_2-s_1).
\]
Thus $\Psi_1\gamma$ is $L$-Lipschitz.  Since
$\gamma=(\Psi_1\gamma)\circ\sigma$ and $\sigma$ is a nondecreasing
surjection, comparison of partitions gives
$\ell(\gamma)\le\ell(\Psi_1\gamma)$.  The Lipschitz estimate gives the
reverse inequality.  Consequently $\Psi_1\gamma$ has length $L$ and speed
$L$ almost everywhere, and its energy is $L^2$.

For $u<1$, change of variables by $\phi_u$ gives
\[
 E(\Psi_u\gamma)
 =\int_0^1\frac{v(t)^2}{(1-u)+uv(t)/L}\,\dd t.
\]
On $\{v>0\}$, convexity of the reciprocal bounds the integrand by
$(1-u)v^2+uLv$; on $\{v=0\}$ it is zero.  This proves
\eqref{eq:shortening-energy}, including $u=1$ by the preceding argument.
The same change of variables preserves length for $u<1$.  Endpoints are
unchanged, and constant-speed paths have $\sigma(t)=t$.

We verify joint strong continuity, including the possibly flat fibers at
$u=1$.  Fix a smooth embedding of $M$ in a Euclidean space and extend its
metric to a smooth positive definite matrix on a neighborhood of $M$.
Suppose $\gamma_i\to\gamma$ strongly in $H^1$ and $u_i\to u$.
The one-dimensional Sobolev embedding gives uniform convergence of the
paths.  Smoothness of the metric gives $v_i\to v$ in $L^2$, and hence
$L_i\to L$.

If $L>0$, the maps $\sigma_i$ and $\phi_{u_i}$ converge uniformly to
$\sigma$ and $\phi_u$.  Given $s_i\to s$, choose $t_i$ in the inverse
fiber defining $\Psi_{u_i}\gamma_i(s_i)$.  Every subsequential limit $t$
satisfies $\phi_u(t)=s$.  The path $\gamma$ has one value on that fiber:
for $u<1$ the fiber is a singleton, and for $u=1$ this was proved above.
Therefore
\[
 \Psi_{u_i}\gamma_i(s_i)=\gamma_i(t_i)
 \longrightarrow\Psi_u\gamma(s).
\]
A compactness argument makes this convergence uniform in $s$.  If $L=0$,
the image of each output is contained in the image of its input, so uniform
convergence to the constant path follows directly.

The energies also converge.  If $L>0$ and $u<1$, the reciprocals in the
integral formula are uniformly bounded and converge in measure; combining
this with $v_i^2\to v^2$ in $L^1$ proves convergence of the integrals.
If $u=1$, Cauchy--Schwarz and \eqref{eq:shortening-energy} give
\[
 L_i^2\le E(\Psi_{u_i}\gamma_i)
 \le(1-u_i)E(\gamma_i)+u_iL_i^2\longrightarrow L^2.
\]
If $L=0$, the upper bound $E(\Psi_{u_i}\gamma_i)\le E(\gamma_i)$ tends
to zero.

For completeness, energy convergence here implies strong, rather than
merely weak, convergence of derivatives.  Uniform convergence and the
energy bound give weak $H^1$ convergence to the already identified output.
In the squared derivative norms, replace the metric matrix at the varying
outputs by its value at the limiting output.  The resulting error tends
to zero, because the matrices converge uniformly and the derivatives are
bounded in $L^2$.  The derivative norms thus converge in one fixed positive
weighted $L^2$ Hilbert norm.  Weak convergence together with convergence of
these norms gives strong convergence.  This proves the asserted joint
continuity.

Finally, for the reversed path the arclength map is
$1-\sigma(1-t)$.  This identity proves commutation with reversal, also
at the quotient endpoint $u=1$.
\end{proof}

We shall use this homotopy only for ordinary homology.  Fixed-time
concatenation of finitely many $H^1$ paths is continuous and adds their
lengths; Lemma~\ref{lem:arclength} then puts the concatenation at constant
speed.  The same is true when some paths are constant.  Restriction to a
moving initial interval, followed by linear rescaling, is also continuous
in $H^1$, including a collapsed interval: for $\gamma_t(s)=\gamma(ts)$,
\[
 E(\gamma_t)=t\int_0^t|\dot\gamma(r)|^2\,\dd r\longrightarrow0
 \quad\text{as }t\longrightarrow0.
\]
Continuity away from zero follows from continuity of dilation in $L^2$.

\subsection{A bounded local path replacement}

The following elementary construction is the source of the uniform
constant in the shortening argument.

\begin{lemma}[Local replacement of paths]\label{lem:shortening-recipes}
Fix $L_0<\infty$ and let $M$ be compact and simply connected.  Pairs of
paths with common endpoints and Lipschitz constants at most $L_0$ admit a
finite cover by relative $C^0$ neighborhoods in this path-pair space
with the following property.  On each
neighborhood there is a relative-endpoint homotopy from the first path to
the second, continuous in the strong $H^1$ topology of the input paths and
the homotopy parameter.  Every path occurring in these homotopies has
length at most a constant $K=K(M,L_0)$.
\end{lemma}

\begin{proof}
The space of such pairs is compact in $C^0$ by the Arzel\`a--Ascoli
theorem.  Consider one pair $(\alpha_0,\beta_0)$ from this space, with
endpoints $x_0,y_0$.  Simple connectivity gives a continuous homotopy
between them, fixing the endpoints.  Replace its paths by polygons on one
sufficiently fine, fixed subdivision of the path parameter.  Uniform
continuity of the original homotopy ensures that all consecutive vertices
are within the injectivity radius.  Polygonal interpolation then gives an
$H^1$-continuous homotopy with a finite length bound.  Short pointwise
geodesic interpolation joins its two endpoint polygons to
$\alpha_0,\beta_0$.  These interpolations also have a finite length
bound, since the original endpoint paths are Lipschitz.

This construction extends to nearby pairs with varying common endpoints
$x,y$.  Attach the short geodesics from $x$ to $x_0$ and from $y_0$ to $y$
to each prototype path, using fixed, small positive time intervals at its
ends.  First choose those intervals small enough that the retimed
prototype endpoints are uniformly close to the original endpoints; this
uses their Lipschitz bounds.  Then restrict $x,y$ and the input paths to
a sufficiently small $C^0$ neighborhood.  Short pointwise geodesic
interpolation joins the actual first path to the adjusted prototype first
path, and similarly at the other end.  Both interpolations fix $x,y$.
Their derivatives are bounded by a fixed multiple of the derivatives of
the two interpolated paths.  Consequently their lengths have a common
bound on this neighborhood.  Every operation just described is continuous
in the strong $H^1$ topology.  A finite subcover of the compact pair space,
and the maximum of its finitely many bounds, give $K$.
\end{proof}

Here and below a local homotopy is a specified choice on a neighborhood;
no global choice of a minimizing geodesic is needed.  We next use these
local choices to shorten a long path block with bounded overhead.

\begin{lemma}[Shortening one block]\label{lem:shortening-block}
Let $M$ be compact and simply connected.
There are constants $C,D>0$, depending only on $M$, with the following
property.  Suppose a path is divided into finitely many pieces, each of
which, on a unit interval, has Lipschitz constant at most one.  Near any
such path, within this space of piecewise Lipschitz paths, there is a
relative-endpoint homotopy to a path of length at
most $D$.  At every stage its length is at most its original length plus
$C$.  The homotopy is continuous in $H^1$ on that neighborhood.  The
constants are independent of the number of pieces.
\end{lemma}

\begin{proof}
Set $D=\operatorname{diam}(M)+1$.  For a fixed prototype choose shortest
paths, called spokes, from its initial point to all its piece vertices.
For nearby inputs extend these choices by short endpoint connectors,
using fixed thirds of the unit interval.  Shrink the neighborhood so that
each spoke has length at most $D$ and Lipschitz constant at most $3D$.
At the initial vertex use the constant spoke.

Proceed through the pieces in order.  Assume the preceding prefix has
already been replaced by its spoke.  Retime that spoke and the next
piece to halves of the enlarged prefix.  This retiming preserves length
and is continuous in $H^1$: every nonconstant path has a positive time
interval throughout the retiming.  In the first step, where the prefix
is empty, insert a constant wait at the initial point and retime the
first piece.  The two paths to be compared are now the spoke followed by
the next piece, and the spoke to the next vertex.  On a unit interval
both have Lipschitz constant at most
\[
 L_0=2\max\{1,3D\}.
\]
Apply Lemma~\ref{lem:shortening-recipes} to replace the first by the
second, and leave the remaining suffix unchanged.

During a replacement the active prefix has length at most $K(M,L_0)$;
during a retiming it has length at most $D+1$.  The suffix still consists
of original pieces.  Thus $C=\max\{K(M,L_0),D+1\}$ bounds the increase
over the length of the entire original path at every stage.  At the end
the path is its final spoke, of length at most $D$.
For a particular input there are finitely many pieces, hence finitely
many local choices.  Intersecting their neighborhoods makes the entire
construction continuous on one neighborhood.  This affects the
neighborhood, but neither $C$ nor $D$.
\end{proof}

The intermediate energies in Lemma~\ref{lem:shortening-block} need not
have a uniform bound: a long prefix may be assigned a short time interval.
Only its length estimate is used.  All time intervals for nonconstant
pieces remain positive, and the construction is an $H^1$ homotopy.
Lemma~\ref{lem:arclength} restores constant speed whenever required.

\subsection{Coloring a finite family}

\begin{lemma}[A fixed-dimensional family shortens by one]
\label{lem:shortening-family}
Let $M$ be compact and simply connected.
For each integer $j\ge0$ there are $a_j,H_j>0$ such that every
$H^1$-continuous family $\Gamma:P\to\mathcal A_a$, where $P$ is a finite
polyhedron of dimension at most $j$ and $a\ge a_j$, is homotopic to a
family in $\mathcal A_{a-1}$ through $\mathcal A_{a+H_j}$.
The constants do not depend on $P$ or on the number of its simplices.
\end{lemma}

\begin{proof}
Normalize the family by Lemma~\ref{lem:arclength}.  A loop of length
$L<a$ now has constant speed $L$.  Let $C,D$ be the constants of
Lemma~\ref{lem:shortening-block}.  Choose $B$ so large that
\begin{equation}\label{eq:shortening-B}
 B/2>D+jC+2.
\end{equation}
For sufficiently large $a$, the integer $k=\lfloor a/B\rfloor$ is at
least $j+1$.  Divide the loop into $k$ equal time blocks, and each block
into $\lceil a/k\rceil$ equal pieces.  Every piece, rescaled to a unit
interval, has Lipschitz constant at most one.  Every block has length
$L/k$.

For each parameter point, Lemma~\ref{lem:shortening-block} supplies
local recipes for the first $j+1$ blocks.  There are only finitely many
blocks and pieces in this particular family.  Choose a common
neighborhood for these recipes, and cover $P$ by finitely many such
neighborhoods.  Subdivide a triangulation finely enough that the closed
stars in a final barycentric subdivision are subordinate to this cover.
Color each vertex by the dimension of the simplex of the preceding
triangulation whose barycenter it is.  The colors lie in
$\{0,\ldots,j\}$ and are distinct on each simplex.

At a vertex $v$ of color $h$, use its local recipe on block $h+1$.
If $\lambda_v$ is its barycentric coordinate, at homotopy time $s$ run
that recipe for the fraction
\[
 s\min\{1,(j+1)\lambda_v\}.
\]
Two positive-coordinate stars of the same color are disjoint; at the
boundary of either star the fraction is zero and the recipe is the
identity.  Recipes of different colors affect different blocks, keeping
their endpoints fixed.  They therefore paste to one continuous $H^1$
homotopy.  At each parameter at most $j+1$ blocks are changed, and the
length throughout is less than $a+(j+1)C$.

At the terminal time some barycentric coordinate is at least
$1/(j+1)$, so one block has been completely shortened.  The final length
is consequently at most
\begin{equation}\label{eq:shortening-terminal}
 L-L/k+D+jC.
\end{equation}
If $L\ge a/2$, then $k\le a/B$ gives $L/k\ge B/2$, and
\eqref{eq:shortening-B} makes \eqref{eq:shortening-terminal} less than
$a-1$.  If $L<a/2$, the rough bound
$L+(j+1)C<a/2+(j+1)C$ has the same consequence once
$a>2((j+1)C+1)$.  Choose $a_j$ larger than this threshold and large
enough that $k\ge j+1$.  We may take $H_j=(j+1)C$.
\end{proof}

\subsection{Filling the remaining bounded family}

\begin{lemma}[A fixed bounded level]\label{lem:shortening-bounded}
For every $b>0$ and $j\ge0$ there is $R_{b,j}<\infty$ such that every
finite $j$-cycle in $\mathcal A_b$ which is null-homologous in $\Lambda$
bounds a finite chain in $\mathcal A_{R_{b,j}}$.
\end{lemma}

\begin{proof}
Normalize by Lemma~\ref{lem:arclength}; the resulting cycle has
$F=\ell<b$.  We describe a single finite model for all such normalized
cycles.  Choose an integer $N$ large enough that paths of energy at most
$(b+1)^2$ move less than a fixed small injectivity radius on intervals
of length $1/N$.  Sampling at these times and joining successive
vertices by minimizing geodesics gives a continuous polygonal map.
For the discrete energy
\[
 E_N(x_0,\ldots,x_{N-1})
   =N\sum_{i=0}^{N-1}d(x_i,x_{i+1})^2,
 \qquad x_N=x_0,
\]
Cauchy--Schwarz gives $E_N\le E$ on sampled loops.
Choose $b'\in(b,b+1)$ such that $(b')^2$ is a regular value of $E_N$.
The set
\[
 Q=\{E_N\le(b')^2\}
\]
is contained strictly inside the open set of vertex tuples whose
successive distances are smaller than the injectivity radius.  It is a
compact smooth manifold with boundary and therefore has finite-dimensional
homology in every degree.  Polygonal interpolation maps $Q$ to loops
of length at most $b'$.

There is also a homotopy from each normalized loop to its sampled
polygon, with one uniform energy bound depending only on $b$ and $M$.
Indeed, on each interval both paths stay in a small normal neighborhood
of its initial vertex.  Pointwise short-geodesic interpolation is
$H^1$-continuous, and the derivative of this interpolation is bounded
by a fixed multiple of the two path derivatives.  Its square-root energy
is thus at most $2C_0b$ for a fixed metric constant $C_0$.  This also
gives a uniform length bound for the homotopy.

The kernel of $H_j(Q;\kk)\to H_j(\Lambda;\kk)$ is finite-dimensional.
Choose finite cycles representing a basis and finite global fillings of
their polygonal images.  Their supports are compact, so their lengths
have a common finite bound.  The sampled image of any cycle in the
lemma represents a linear combination of this basis.  The difference
bounds in $Q$, whose entire image has length at most $b'$.  Adding that
comparison chain, the finitely many chosen fillings, and the two
normalization and polygonal homotopies proves the claim with one
$R_{b,j}$ strictly larger than all these bounds.  No bound on the number
of simplices in the original cycle is needed.
\end{proof}

\begin{proposition}[Fixed-degree filling]\label{prop:fixed-degree-death}
Let $M$ be compact and simply connected.  For every integer $j\ge0$
there is $D_j>0$ such that, for every real $a>0$, each finite $j$-cycle
in $\mathcal A_a$ which is null-homologous in $\Lambda$ bounds a finite
chain in $\mathcal A_{a+D_j}$.
\end{proposition}

\begin{proof}
A finite singular cycle over $\kk$ can be regarded as the image of a
cycle on a finite polyhedron of dimension at most $j$.  One explicit
construction takes the finitely many singular simplices, pairs equal
faces that cancel in their boundary, and glues these faces by their
ordered affine identifications.  The first barycentric subdivision has
embedded closed simplices and vertices distinguished by their original
face dimensions; a second subdivision removes multiple simplices with
the same vertex set.  This gives a finite polyhedron with its
top-dimensional cycle and the original map to $\Lambda$.
The usual subdivision and prism chains identify its image with the
original singular cycle.  All these chains have image in the union of
the original simplices.  Thus the number of parameter simplices plays
no role in a length estimate.

Apply Lemma~\ref{lem:shortening-family} repeatedly, normalizing between
successive applications.  Starting at level $a$, the successive
starting levels are $a,a-1,a-2,\ldots$, until the family lies below one
fixed level $b_j>a_j+1$.  The prism chains at the successive stages
all lie in $\mathcal A_{a+H_j}$: the peak at stage $r$ is bounded by
$a-r+H_j$, so these bounds are not added over the stages.  Only finitely
many stages are needed for each input cycle.

The resulting cycle is still globally null-homologous.  Apply
Lemma~\ref{lem:shortening-bounded} at the fixed level $b_j$.  Taking
$D_j>\max\{H_j,R_{b_j,j}\}$ gives the required filling after adjoining
the shortening prisms.  If the original $a$ is already at most $b_j$,
the bounded-level lemma applies immediately.  This proves the assertion
for every positive level, without a regular-value assumption.
\end{proof}

The constants $D_j$ presently depend on the degree.  For spheres, the
cut--join operations in the next section reduce all sufficiently large
degrees to a fixed finite range.  This will turn
Proposition~\ref{prop:fixed-degree-death} into one filling bound valid
in every degree.
